% from ega4-16-fr.tex
\oldpage[IV-4]{26}
Le diagramme commutatif
\[
  \xymatrix{
    Y\ar[d]_-g&X\times_S Y\ar[l]_-q\ar[d]_-h&
      X\times_S Y\ar[l]_-{\operatorname{id}}\ar[d]^-p\\
    S&S\ar[l]^-{\operatorname{id}}&X\ar[l]^-f
  }
\]
donne une factorisation de l'\emph{isomorphisme} canonique $P^n(p)$ (16.4.5)
\[
  p^*(\mathscr{P}_{X/S}^n)\longrightarrow
  \mathscr{P}_{Z/S}^n\longrightarrow\mathscr{P}_{Z/Y}^n
\]
et de même, en intervertissant les rôles de $X$ et $Y$, on a une factorisation
de l'\emph{isomorphisme} $P^n(q)$
\[
  q^*(\mathscr{P}_{Y/S}^n)\longrightarrow
  \mathscr{P}_{Z/S}^n\longrightarrow\mathscr{P}_{Z/X}^n.
\]
Ceci prouve que l'homomorphisme canonique (16.4.18.1)
\[
  p_{Z/X/S}:p^*(\mathscr{P}_{X/S}^n)\longrightarrow\mathscr{P}_{Z/S}^n
  \quad\text{(resp.\
  $q_{Z/Y/S}:q^*(\mathscr{P}_{Y/S}^n)\to\mathscr{P}_{Z/S}^n$)}
\]
est \emph{injectif}, et que le noyau de l'homomorphisme canonique surjectif
(16.4.18.2)
\[
  \mathscr{P}_{Z/S}^n\longrightarrow\mathscr{P}_{Z/Y}^n
  \quad\text{(resp.\
  $\mathscr{P}_{Z/S}^n\to\mathscr{P}_{Z/X}^n$)}
\]
est supplémentaire de l'image de $p_{Z/X/S}$ (resp.\ $q_{Z/Y/S}$). Mais
d'autre part, ce noyau est, en vertu de (16.4.18), engendré par l'image par
$q_{Z/Y/S}$ (resp.\ $p_{Z/X/S}$) de l'Idéal d'augmentation de
$q^*(\mathscr{P}_{Y/S}^n)$ (resp.\ $p^*(\mathscr{P}_{X/S}^n)$). On en conclut
la proposition en considérant le cas $n=1$.

On généralise aussitôt (16.4.23) au cas d'un produit d'un nombre fini
quelconque de $S$-préschémas.

\begin{remarks}[16.4.24]
\label{IV.16.4.24}
\begin{enumerate}
\item[(i)] Nous verrons (17.2.3) que lorsque le morphisme $f:X\to Y$ dans
(16.4.18) est \emph{lisse}, l'homomorphisme $f_{X/Y/Z}$ dans (16.4.19.1) est
localement \emph{inversible à gauche} et en particulier injectif. De même,
lorsque le morphisme $f\circ j:X\to Z$ de (16.4.20) est \emph{lisse},
l'homomorphisme de gauche dans (16.4.21.2) est localement
\emph{inversible à gauche} et \emph{a fortiori} injectif (17.2.5). Au chap. V,
nous donnerons aussi une variante, dans le cas des Modules sur les
préschémas, des «modules d'imperfection» étudiés dans (0, 20.6), et des suites
exactes où ils figurent.

\item[(ii)] Soient $X$ un espace topologique, $\mathscr{A}$ un faisceau
d'anneaux sur $X$ et $\mathscr{B}$ une $\mathscr{A}$-Algèbre sur $X$. Alors il
est clair que
\[
  U\longmapsto
  P_{\Gamma(U,\mathscr{B})/\Gamma(U,\mathscr{A})}^n
  \qquad (U\text{ ouvert dans }X)
\]
est un préfaisceau de $\Gamma(U,\mathscr{B})$-algèbres augmentées, donc le
faisceau associé $\mathscr{P}_{\mathscr{B}/\mathscr{A}}^n$ est une
$\mathscr{B}$-Algèbre augmentée. Dans le cas particulier où $X$ est un
préschéma, $f=(\psi,\theta):X\to S$ un morphisme de préschémas, il résulte
aisément de (16.4.16) et de l'exactitude du foncteur $\varinjlim$ que
$\mathscr{P}_{X/S}^n$ est canoniquement isomorphe à
$\mathscr{P}_{\mathcal{O}_X/\psi^*(\mathcal{O}_S)}^n$. Il s'ensuit que le
formalisme développé dans le présent paragraphe pourrait être considéré comme
un
